\documentclass[conference]{IEEEtran}
\IEEEoverridecommandlockouts

\usepackage{cite}
\usepackage{amsmath,amssymb,amsfonts}
\usepackage{amsthm}
\usepackage{algorithmic}
\usepackage{algorithm}
\usepackage{graphicx}
\usepackage{textcomp}
\usepackage{xcolor}
\usepackage{booktabs}
\usepackage{multirow}
\usepackage{url}
\usepackage{bm}

\newtheorem{definition}{Definition}
\newtheorem{proposition}{Proposition}

\def\BibTeX{{\rm B\kern-.05em{\sc i\kern-.025em b}\kern-.08em
    T\kern-.1667em\lower.7ex\hbox{E}\kern-.125emX}}

\begin{document}

\title{SW-LTS: Similarity-Weighted Linguistic Time Series Forecasting\\
via Gaussian Kernel Aggregation over Hedge Algebra Semantic Space}

\author{
\IEEEauthorblockN{Nguyen Duy Hieu}
\IEEEauthorblockA{
\textit{Faculty of Natural Sciences}\\
Tay Bac University\\
Son La, Viet Nam\\
hieund@utb.edu.vn}
}

\maketitle

% ─────────────────────────────────────────────────────────────────────────────
\begin{abstract}
Linguistic Time Series (LTS) forecasting, grounded in the algebraic
theory of Hedge Algebras (HA), maps numeric observations to linguistic
labels and derives forecasting rules through Linguistic Logical
Relationship Groups (LLRGs). A fundamental limitation shared by all
existing LTS variants is their reliance on \emph{exact-match lookup}:
forecasting succeeds only when the current linguistic pattern has been
observed verbatim during training, discarding all information from
semantically similar but non-identical rules. We formalize this as the
\emph{sparse-rule problem}: as the forecasting order $\lambda$ increases,
the space of possible Left-Hand Side (LHS) patterns grows exponentially
as $|W|^\lambda$, while training observations remain fixed. We propose
\textbf{SW-LTS} (Similarity-Weighted LTS), which replaces exact-match
lookup with a Gaussian kernel-weighted aggregation over the full LLRG.
The kernel operates in normalized HA semantic space, yielding a
dimensionless bandwidth parameter $\sigma$ that is scale-invariant
across datasets. Theoretical analysis establishes that SW-LTS strictly
subsumes both LTS ($\sigma \to 0$) and the global-mean fallback of
High-Order LTS (HO-LTS, $\sigma \to \infty$). Experiments on four
standard benchmarks—Alabama Enrollment, TAIEX 1999, Wolf's Sunspot,
and Tuscaloosa Temperature—show that SW-LTS achieves the lowest or
joint-lowest error across all metrics, reducing Mean Squared Error
(MSE) by 4.4\% over the strongest baseline on Alabama. The method
requires no architectural changes beyond the baseline LTS pipeline.
\end{abstract}

\begin{IEEEkeywords}
fuzzy time series, hedge algebras, linguistic time series,
kernel-weighted forecasting, sparse-rule problem, similarity-based rule aggregation
\end{IEEEkeywords}

% ─────────────────────────────────────────────────────────────────────────────
\section{Introduction}
\label{sec:intro}

Time series forecasting is fundamental to applications in finance,
meteorology, energy management, and public health. Classical parametric
methods such as ARIMA require stationarity assumptions, while modern
deep learning approaches demand large corpora and offer limited
interpretability~\cite{zhan2024dfcnn,wang2025bn}. Fuzzy Time Series
(FTS), introduced by Song and Chissom~\cite{song1993,song1994},
addresses these constraints by encoding observations as linguistic
labels and deriving forecasting rules from fuzzy logical relationships.
The resulting models are transparent, computationally light, and
effective on short series. A broad family of FTS extensions has since
emerged, refining the partition strategy, rule defuzzification, and
weighting schemes~\cite{chen1996,yu2005}.

Linguistic Time Series (LTS)~\cite{hieu2020lts} advances FTS by
grounding the linguistic encoding in the formal algebraic theory of
\textit{Hedge Algebras} (HA)~\cite{catho2002ha,catho2002sqm}. Unlike
equal-width interval partitioning, HA generates a structured vocabulary
of linguistic terms from first principles, assigns each term a
Semantic Quantifying Measure (SQM) value derived axiomatically from
two shape parameters $(\theta, \alpha)$, and embeds all terms in a
total semantic order. This endows LTS with a rigorous semantic
foundation. Subsequent extensions added high-order rule extraction
(High-Order LTS, HO-LTS~\cite{hieu2021holts}), Particle Swarm
Optimization (PSO)-based parameter tuning
(LTS-PSO~\cite{hieu2022ltspso}), and nested co-optimization of both
parameters and vocabulary
(Co-Optimization LTS, CO-LTS~\cite{hieu2023colts}).

Despite these advances, all LTS variants share a structural limitation:
\emph{exact-match lookup}—the process of querying the LLRG with the
current order-$\lambda$ linguistic pattern and producing a forecast
\emph{only} when that exact pattern was observed during training. When
the lookup fails, the model reverts to a coarse heuristic, discarding
all information contained in semantically adjacent rules. We call this
the \textbf{sparse-rule problem}: as $\lambda$ grows, the number of
possible Left-Hand Side (LHS) patterns increases as $|W|^\lambda$ while
the training set size stays fixed, leaving an exponentially growing
fraction of patterns unmatched. Although the phenomenon of unmatched
states has been implicitly acknowledged in the FTS literature—addressed
through interval partitioning adjustments and mean-value fallbacks
(e.g., HO-LTS~\cite{hieu2021holts})—no prior work has formally
characterized it as a function of order and vocabulary size, nor
proposed a principled kernel-based remedy rooted in the HA semantic
metric.

We propose \textbf{SW-LTS} (Similarity-Weighted LTS), which replaces
binary lookup with continuous Gaussian kernel-weighted aggregation
over all LLRG entries. The kernel is computed in normalized HA semantic
space, so its bandwidth $\sigma$ is dimensionless and dataset-agnostic.

\textbf{Contributions:}
\begin{enumerate}
  \item A kernel-weighted forecasting mechanism that exploits the full
        LLRG for every query, eliminating the binary matched/unmatched
        dichotomy.
  \item Formal proofs that SW-LTS recovers LTS as $\sigma \to 0$ and
        the HO-LTS global-mean fallback as $\sigma \to \infty$,
        establishing it as a strict generalization of both.
  \item A scale-invariant $\sigma$ formulation that enables consistent
        grid search across datasets with disparate numerical ranges.
  \item Empirical validation—including a rule coverage analysis—on four
        standard benchmarks against five representative baselines.
\end{enumerate}

Section~\ref{sec:background} reviews HA and LTS.
Section~\ref{sec:method} presents SW-LTS.
Section~\ref{sec:exp} reports experiments.
Section~\ref{sec:conclusion} concludes.

% ─────────────────────────────────────────────────────────────────────────────
\section{Background and Motivation}
\label{sec:background}

\subsection{Hedge Algebras}

Hedge Algebras (HA)~\cite{catho2002ha} provide an algebraic framework
for linguistic variables. An HA is a quintuple
$\mathcal{AX} = (X, G, C, H, \leq)$, where:
\begin{itemize}
  \item $X$ is a set of linguistic terms (\emph{words}),
  \item $G = \{c^-, c^+\}$ are two generators representing the extreme
        semantic poles (\emph{Small} and \emph{Large}),
  \item $C = \{w_0\}$ is a semantically neutral element,
  \item $H$ is a set of hedge operators (linguistic modifiers), and
  \item $\leq$ is a linear order on $X$ consistent with the semantic
        ordering induced by $G$ and $H$.
\end{itemize}

Two hedges are of primary interest: $V$ (\emph{Very}) strengthens a
term while $L$ (\emph{Rather}/\emph{Little}) weakens it. Applying
these hedges recursively to $c^-$ and $c^+$ generates a hierarchy of
words. At specificity level $k$ the vocabulary contains $2^{k+2}-1$
words (level $k=1$: 7 words; $k=2$: 15 words; $k=3$: 31 words).

The \textbf{Semantic Quantifying Measure} (SQM)~\cite{catho2002sqm},
$\nu : X \to [0, 1]$, assigns each word a unique value that
monotonically respects $\leq$. The SQM is fully determined by two
parameters: $\theta = \nu(c^-)$ and $\alpha$, the weight of the hedge
$L$. Given $(\theta, \alpha)$, all word SQM values are computed
analytically. The \textbf{semantic point} of a word $w$ with respect
to the universe of discourse $[\ell, u]$ is then
\begin{equation}
  s(w) = \ell + (u - \ell)\,\nu(w).
\label{eq:semantic_point}
\end{equation}

\subsection{Linguistic Time Series Model}

The LTS model~\cite{hieu2020lts} forecasts a time series
$\{y_t\}_{t=1}^{T}$ via a six-step pipeline.

\textbf{Step 1 — Vocabulary selection.}
Choose a specificity level $k$ to obtain the word set $W$ of size
$|W| = 2^{k+2}-1$.

\textbf{Step 2 — Semantic quantification.}
Compute $\nu(w)$ for every $w \in W$ using $(\theta, \alpha)$.

\textbf{Step 3 — Semantic points.}
Map each $\nu(w)$ to the universe $[\ell, u]$ via \eqref{eq:semantic_point}.

\textbf{Step 4 — Fuzzification.}
Assign each observation the nearest word by semantic distance:
$\ell_t = \arg\min_{w \in W} |s(w) - y_t|$.

\textbf{Step 5 — Rule extraction.}
For $t = \lambda{+}1, \ldots, T$, record the order-$\lambda$
Linguistic Logical Relationship (LLR):
$(\ell_{t-\lambda}, \ldots, \ell_{t-1}) \to \ell_t$.
All LLRs sharing the same LHS tuple are aggregated into a Linguistic
Logical Relationship Group (LLRG):
$\mathcal{R} : (w_1,\ldots,w_\lambda) \mapsto [r_1, r_2, \ldots]$.

\textbf{Step 6 — Forecasting.}
For the query window
$\bm{\ell}_t = (\ell_{t-\lambda}, \ldots, \ell_{t-1})$:
\begin{equation}
  \hat{y}_t =
  \begin{cases}
    \dfrac{1}{|\mathcal{R}(\bm{\ell}_t)|}
    \displaystyle\sum_{r \in \mathcal{R}(\bm{\ell}_t)} s(r)
    & \text{if } \bm{\ell}_t \in \mathcal{R}, \\[6pt]
    s(\ell_{t-1})
    & \text{otherwise (no-match fallback).}
  \end{cases}
\label{eq:lts_forecast}
\end{equation}

\subsection{Limitations of Existing LTS Models}

\textbf{L1 — Sparse-rule problem.}
Equation~\eqref{eq:lts_forecast} applies the same uninformative
fallback $s(\ell_{t-1})$ to \emph{every} unmatched LHS pattern,
regardless of semantic proximity to existing LLRG entries. The rule
coverage ratio $|\mathcal{R}|/|W|^\lambda$ decreases exponentially
with $\lambda$: at $\lambda=3$ with $|W|=15$, the pattern space
contains $3{,}375$ tuples, yet a small dataset may populate fewer
than 20 of them (coverage $< 1\%$).

\textbf{L2 — Discontinuous forecast surface.}
LTS produces step-function forecasts over the LLRG key space: two
query windows differing by a single hedge application receive entirely
different forecasts if they fall into different LLRG entries, even
when their semantic distance is negligible.

\textbf{L3 — No graded confidence.}
A query that closely matches an LLRG entry and one with no nearby
entry are treated identically, providing no mechanism to reflect
varying reliability of the forecast.

HO-LTS~\cite{hieu2021holts} partially mitigates L1 by replacing the
no-match fallback with the mean semantic point of all words in the LHS
tuple, but retains exact-match lookup for the matched case.
LTS-PSO~\cite{hieu2022ltspso} and CO-LTS~\cite{hieu2023colts} optimize
$(\theta, \alpha)$ without modifying the lookup mechanism. SW-LTS
addresses all three limitations simultaneously.

% ─────────────────────────────────────────────────────────────────────────────
\section{The Proposed SW-LTS Model}
\label{sec:method}

\subsection{Kernel-Weighted Forecasting Framework}

SW-LTS replaces the binary LLRG lookup with a continuous
kernel-weighted average over \emph{all} LLRG entries. Semantic
similarity between the query window and each rule's LHS is measured
by a Gaussian (Radial Basis Function, RBF) kernel in normalized HA
semantic space.

\begin{definition}[LHS Semantic Vector]
For a pattern $\bm{w} = (w_1, \ldots, w_\lambda) \in W^\lambda$, the
\emph{semantic vector} is
$\phi(\bm{w}) = \bigl(s(w_1),\, \ldots,\, s(w_\lambda)\bigr) \in \mathbb{R}^\lambda.$
\end{definition}

\begin{definition}[Normalized Semantic Vector]
Given $[\ell, u]$ with span $\Delta = u - \ell > 0$, the
\emph{normalized semantic vector} is
\[
  \tilde{\phi}(\bm{w})
  = \frac{1}{\Delta}\bigl(s(w_1)-\ell,\, \ldots,\, s(w_\lambda)-\ell\bigr)
  \;\in\; [0,1]^\lambda.
\]
\end{definition}

Because all semantic points satisfy $s(w) \in [\ell, u]$, the
normalized vector is confined to the unit hypercube $[0,1]^\lambda$
regardless of the data scale, which is the key to scale-invariance.

\subsection{Gaussian Kernel Similarity}

The similarity between two patterns $\bm{u}$ and $\bm{v}$ is
\begin{equation}
  K(\bm{u}, \bm{v})
  = \exp\!\left(-\frac{\|\tilde{\phi}(\bm{u}) - \tilde{\phi}(\bm{v})\|^2}
                      {2\sigma^2}\right),
\label{eq:kernel}
\end{equation}
where $\sigma > 0$ is the bandwidth parameter. Since $\tilde{\phi}$
maps to $[0,1]^\lambda$, $\sigma$ is dimensionless: $\sigma = 0.05$
means the effective kernel radius covers 5\% of the normalized range
per dimension, regardless of whether the raw data span is tens or
thousands of units.

\subsection{SW-LTS Forecasting Formula}

Let $\mathcal{R} = \{(\bm{k}_j, R_j)\}_{j=1}^{M}$ denote the LLRG,
where $\bm{k}_j$ is the LHS pattern and $R_j$ the set of observed
Right-Hand Side (RHS) labels of the $j$-th group. The SW-LTS forecast
for query window $\bm{\ell}_t$ is
\begin{equation}
  \hat{y}_t^{\mathrm{SW}}
  = \frac{\displaystyle\sum_{j=1}^{M}
          K(\bm{\ell}_t, \bm{k}_j)\;\bar{s}(R_j)}
         {\displaystyle\sum_{j=1}^{M}
          K(\bm{\ell}_t, \bm{k}_j)},
\label{eq:sw_forecast}
\end{equation}
where $\bar{s}(R_j) = \frac{1}{|R_j|}\sum_{r \in R_j} s(r)$ is the
mean semantic point of the RHS labels of group $j$. If
$\mathcal{R} = \emptyset$ or the denominator is numerically zero, the
fallback $\hat{y}_t = s(\ell_{t-1})$ is applied.

\subsection{Theoretical Properties}

\begin{proposition}[$\sigma \to 0$: Recovery of LTS exact-match]
As $\sigma \to 0$, $K(\bm{\ell}_t, \bm{k}_j) \to 1$ if
$\bm{\ell}_t = \bm{k}_j$ and $\to 0$ otherwise. If
$\bm{\ell}_t \in \mathcal{R}$, all weight concentrates on the
matching group and \eqref{eq:sw_forecast} reduces to
\eqref{eq:lts_forecast}. If $\bm{\ell}_t \notin \mathcal{R}$, all
weights vanish and the fallback applies. SW-LTS thus recovers
LTS in the small-$\sigma$ limit.
\end{proposition}

\begin{proposition}[$\sigma \to \infty$: Recovery of HO-LTS fallback]
As $\sigma \to \infty$, $K(\bm{u}, \bm{v}) \to 1$ for all pairs.
Formula \eqref{eq:sw_forecast} becomes the unweighted mean of all
group RHS averages: $\hat{y}_t = \frac{1}{M}\sum_j \bar{s}(R_j)$.
This coincides with the HO-LTS global-mean fallback when no rule
matches~\cite{hieu2021holts}.
\end{proposition}

SW-LTS is therefore a \emph{strict generalization} of both LTS and
the HO-LTS fallback, with $\sigma$ smoothly interpolating between
local exact-match specificity and global semantic averaging.

\subsection{Bandwidth Selection}

In the unit hypercube $[0,1]^\lambda$, a practical data-driven default
for $\sigma$ is the median pairwise distance between normalized
semantic points of words in $W$. For systematic comparison, we use
grid search over $\sigma \in \{0.01, 0.03, 0.05, 0.08, 0.10, 0.15,
0.20, 0.30, 0.50, 0.80\}$ with minimum in-sample MSE as the selection
criterion. Because $\sigma$ is normalized, the same grid is applicable
across all datasets without rescaling.

\subsection{Joint Optimization: SW-PSO-LTS}

To jointly optimize the HA parameters $(\theta, \alpha)$ and the
bandwidth $\sigma$, we further propose \textbf{SW-PSO-LTS}: a PSO
wrapper~\cite{kennedy1995pso} that minimizes in-sample MSE over the
search space $(\theta,\alpha,\sigma) \in [0.3,0.7]^2 \times [0.01,1.0]$
using $N=100$ particles, $G_{\max}=500$ iterations, inertia weight
$\omega=0.4$, and cognitive/social coefficients $c_1=c_2=2.0$. Three
independent runs are executed and the best-MSE solution is retained.

\subsection{Algorithm}

\begin{algorithm}
\caption{SW-LTS: Training and Forecasting}
\begin{algorithmic}[1]
\REQUIRE Series $\{y_t\}_{t=1}^{T}$, universe $[\ell, u]$,
         parameters $\theta$, $\alpha$, $\sigma$,
         order $\lambda$, specificity $k$
\STATE Compute $\nu(w)$ and $s(w)$ for all $w \in W$
\STATE Fuzzify: $\ell_t \leftarrow \arg\min_{w \in W} |s(w) - y_t|$,
       $\forall t$
\STATE Build LLRG $\mathcal{R}$ from order-$\lambda$ LLRs
\STATE Normalize: $\tilde{s}(w) \leftarrow (s(w)-\ell)/(u-\ell)$
       for all $w$
\FOR{$t = \lambda{+}1, \ldots, T$}
  \STATE $\bm{\ell}_t \leftarrow (\ell_{t-\lambda}, \ldots, \ell_{t-1})$
  \STATE $w_j \leftarrow \exp\!\bigl(-\|\tilde{\phi}(\bm{\ell}_t) -
         \tilde{\phi}(\bm{k}_j)\|^2 / 2\sigma^2\bigr)$ for all $j$
  \IF{$\sum_j w_j > 0$}
    \STATE $\hat{y}_t \leftarrow \sum_j w_j \bar{s}(R_j) / \sum_j w_j$
  \ELSE
    \STATE $\hat{y}_t \leftarrow s(\ell_{t-1})$
  \ENDIF
\ENDFOR
\RETURN $\{\hat{y}_t\}_{t=\lambda+1}^{T}$
\end{algorithmic}
\end{algorithm}

% ─────────────────────────────────────────────────────────────────────────────
\section{Experiments}
\label{sec:exp}

\subsection{Datasets}

We use four benchmark datasets standard in the LTS and FTS
literature~\cite{song1993,hieu2020lts,hieu2021holts}:

\begin{table}[htbp]
\centering
\caption{Benchmark Datasets}
\label{tab:datasets}
\begin{tabular}{@{}lrcc@{}}
\toprule
Dataset & $n$ & Universe $[\ell, u]$ & Domain \\
\midrule
Alabama Enrollment & 22  & [13\,000,\; 20\,000] & Education \\
TAIEX 1999         & 241 & [5\,201,\; 9\,039]   & Finance \\
Wolf's Sunspot     & 288 & [0,\; 296.23]         & Astronomy \\
Tuscaloosa Temp.   & 122 & [23,\; 32]            & Meteorology \\
\bottomrule
\end{tabular}
\end{table}

Alabama Enrollment~\cite{song1993} (annual university enrollment,
1971--1992) is the canonical LTS benchmark.
TAIEX 1999~\cite{yu2005} contains daily closing values of the Taiwan
Stock Exchange Capitalization Weighted Stock Index for the year 1999.
Wolf's Sunspot records monthly mean sunspot numbers from 1954 to 1977
(SILSO data, Royal Observatory of Belgium).
Tuscaloosa Temperature gives monthly average temperature in
$^{\circ}$C from a U.S. weather station (NCEI archive).

\subsection{Comparison Methods}

We compare against five baselines spanning the FTS and LTS families:
\begin{itemize}
  \item \textbf{Song \& Chissom (1993)}~\cite{song1993}: original
        FTS with 7 equal-width intervals and max-min composition.
  \item \textbf{Chen (1996)}~\cite{chen1996}: simplified FTS with
        7 equal-width intervals and arithmetic-mean defuzzification.
  \item \textbf{LTS (2020)}~\cite{hieu2020lts}: first-order LTS
        ($\lambda=1$, $k=2$, $\theta=0.57$, $\alpha=0.49$).
  \item \textbf{HO-LTS (2021)}~\cite{hieu2021holts}: high-order
        LTS, evaluated at $\lambda \in \{2,3\}$, $k=2$.
  \item \textbf{LTS-PSO (2022)}~\cite{hieu2022ltspso}: LTS with
        PSO-tuned $(\theta,\alpha)$ and a modified forecasting
        formula $\hat{y}_t = \tfrac{1}{2}(s(\ell_{t-1}) +
        \bar{s}(\mathcal{R}(\bm{\ell}_t)))$ when a rule matches.
\end{itemize}

\subsection{Performance Metrics}

We report Mean Squared Error (MSE), Root Mean Squared Error (RMSE),
Mean Absolute Percentage Error (MAPE), and Mean Absolute Error (MAE):
\begin{align}
  \mathrm{MSE}  &= \tfrac{1}{n}\textstyle\sum_{t}(y_t-\hat{y}_t)^2, \\
  \mathrm{RMSE} &= \sqrt{\mathrm{MSE}}, \\
  \mathrm{MAPE} &= \tfrac{100\%}{n}\textstyle\sum_{t}
    \tfrac{|y_t-\hat{y}_t|}{|y_t|},\quad y_t\neq 0, \\
  \mathrm{MAE}  &= \tfrac{1}{n}\textstyle\sum_{t}|y_t-\hat{y}_t|.
\end{align}
MSE is the primary ranking metric. MAPE is excluded when $y_t = 0$
(relevant for Sunspot, which contains zero-observation months).

\subsection{Experimental Protocol}

All experiments use the \emph{one-step-ahead in-sample} protocol
standard in LTS research~\cite{hieu2020lts,hieu2021holts}: each model
is fitted on the full series, and for each time step $t$, the forecast
$\hat{y}_t$ is produced from the preceding $\lambda$ observations
without look-ahead bias. For LTS-family models, HA parameters are set
to $\theta=0.57$, $\alpha=0.49$ (default from~\cite{hieu2020lts})
unless the method specifies its own optimization procedure. Song \&
Chissom (1993) and Chen (1996) partition $[\ell, u]$ into 7 equal-width
intervals as in their original formulations. All LTS-based methods
use $k=2$ (15-word vocabulary). SW-LTS selects $\sigma$ via the
10-point grid $\{0.01, 0.03, \ldots, 0.80\}$ in $[0,1]^\lambda$ by
minimum in-sample MSE, and is evaluated at $\lambda \in \{2,3\}$.
Experiments were implemented in Python using the open-source
\texttt{pyLTS} library.

\subsection{Quantitative Results}

Tables~\ref{tab:mse}--\ref{tab:mae} compare all methods across four
metrics. \textbf{Bold} marks the best result per column;
\underline{underline} marks the second best.

\begin{table}[htbp]
\centering
\caption{MSE Comparison}
\label{tab:mse}
\begin{tabular}{@{}lrrrr@{}}
\toprule
Method & Alabama & TAIEX-99 & Sunspot & Temp. \\
\midrule
Song \& Chissom (1993) & 806{,}087 & 103{,}899 & 2{,}584 & 3.77 \\
Chen (1996)            & 407{,}521 &  42{,}617 & 3{,}342 & 1.08 \\
LTS (2020)             & 123{,}525 &  26{,}517 & 2{,}363 & 0.79 \\
HO-LTS, $\lambda=2$   &  64{,}082 &  20{,}255 &    907  & 0.44 \\
HO-LTS, $\lambda=3$   &  15{,}077 &  14{,}728 &    562  & 0.16 \\
LTS-PSO (2022)         & 140{,}525 &  15{,}477 & 1{,}382 & 0.77 \\
\midrule
SW-LTS, $\lambda=2$   &  64{,}082 &  18{,}283 &    907  & 0.44 \\
\textbf{SW-LTS}, $\bm{\lambda=3}$ & \textbf{14{,}420} & \underline{14{,}728} & \underline{562} & \underline{0.16} \\
\bottomrule
\end{tabular}
\end{table}

\begin{table}[htbp]
\centering
\caption{RMSE Comparison}
\label{tab:rmse}
\begin{tabular}{@{}lrrrr@{}}
\toprule
Method & Alabama & TAIEX-99 & Sunspot & Temp. \\
\midrule
Song \& Chissom (1993) & 897.82 & 322.33 & 50.83 & 1.94 \\
Chen (1996)            & 638.37 & 206.44 & 57.81 & 1.04 \\
LTS (2020)             & 351.46 & 162.84 & 48.62 & 0.89 \\
HO-LTS, $\lambda=2$   & 253.14 & 142.32 & 30.11 & 0.66 \\
HO-LTS, $\lambda=3$   & 122.79 & 121.36 & 23.70 & 0.40 \\
LTS-PSO (2022)         & 374.87 & 124.41 & 37.18 & 0.88 \\
\midrule
SW-LTS, $\lambda=2$   & 253.14 & 135.21 & 30.11 & 0.66 \\
\textbf{SW-LTS}, $\bm{\lambda=3}$ & \textbf{120.08} & \underline{121.36} & \underline{23.70} & \underline{0.40} \\
\bottomrule
\end{tabular}
\end{table}

\begin{table}[htbp]
\centering
\caption{MAPE (\%) Comparison}
\label{tab:mape}
\begin{tabular}{@{}lrrrr@{}}
\toprule
Method & Alabama & TAIEX-99 & Sunspot$^\dagger$ & Temp. \\
\midrule
Song \& Chissom (1993) & 3.756 & 3.215 &  62.72 & 5.726 \\
Chen (1996)            & 3.110 & 2.357 & 243.90 & 2.996 \\
LTS (2020)             & 1.616 & 1.740 & 206.75 & 2.600 \\
HO-LTS, $\lambda=2$   & 1.168 & 1.444 & 138.39 & 1.785 \\
HO-LTS, $\lambda=3$   & 0.686 & 1.258 & 105.95 & 0.894 \\
LTS-PSO (2022)         & 1.862 & 1.323 & 110.35 & 2.594 \\
\midrule
SW-LTS, $\lambda=2$   & 1.168 & 1.375 & 138.39 & 1.785 \\
\textbf{SW-LTS}, $\bm{\lambda=3}$ & \textbf{0.653} & \underline{1.258} & \underline{105.95} & \underline{0.894} \\
\bottomrule
\end{tabular}\\[3pt]
{\footnotesize $^\dagger$MAPE inflated by near-zero observations; RMSE
and MAE are more reliable for Sunspot.}
\end{table}

\begin{table}[htbp]
\centering
\caption{MAE Comparison}
\label{tab:mae}
\begin{tabular}{@{}lrrrr@{}}
\toprule
Method & Alabama & TAIEX-99 & Sunspot & Temp. \\
\midrule
Song \& Chissom (1993) & 632.05 & 239.92 & 35.52 & 1.66 \\
Chen (1996)            & 498.81 & 172.00 & 51.70 & 0.86 \\
LTS (2020)             & 265.67 & 129.74 & 42.73 & 0.74 \\
HO-LTS, $\lambda=2$   & 191.95 & 107.87 & 22.26 & 0.51 \\
HO-LTS, $\lambda=3$   & 112.38 &  93.13 & 16.18 & 0.25 \\
LTS-PSO (2022)         & 307.33 &  97.19 & 30.98 & 0.74 \\
\midrule
SW-LTS, $\lambda=2$   & 191.95 & 102.37 & 22.26 & 0.51 \\
\textbf{SW-LTS}, $\bm{\lambda=3}$ & \textbf{107.76} & \underline{93.13} & \underline{16.18} & \underline{0.25} \\
\bottomrule
\end{tabular}
\end{table}

\subsection{Rule Coverage Analysis}

The degree to which SW-LTS improves over exact-match LTS is expected
to depend on rule coverage: the fraction of all possible order-$\lambda$
LHS patterns that are populated in the LLRG,
$\mathrm{Cov}(\lambda) = |\mathcal{R}|/|W|^\lambda$.
Table~\ref{tab:coverage} reports coverage for all four datasets at
$\lambda=3$ with $|W|=15$.

\begin{table}[htbp]
\centering
\caption{LLRG Rule Coverage at $\lambda=3$, $k=2$ ($|W|=15$)}
\label{tab:coverage}
\begin{tabular}{@{}lrrrr@{}}
\toprule
Dataset & $n$ & $|W|^\lambda$ & $|\mathcal{R}|$ & Cov.\ (\%) \\
\midrule
Alabama   &  22 & 3{,}375 &  19 & \textbf{0.56} \\
TAIEX-99  & 241 & 3{,}375 &  89 & 2.64 \\
Sunspot   & 288 & 3{,}375 & 139 & 4.12 \\
Temp.     & 122 & 3{,}375 & 105 & 3.11 \\
\bottomrule
\end{tabular}
\end{table}

Alabama's coverage (0.56\%) is 4--7$\times$ lower than the other
three datasets (2.6--4.1\%), which directly predicts the differential
benefit of kernel smoothing observed in the results.

\subsection{Analysis and Discussion}

\textbf{Effect of rule coverage on SW-LTS gain.}
At $\lambda=3$, Alabama's LLRG contains only 19 entries out of 3,375
possible LHS tuples. The exact-match mechanism in LTS and HO-LTS
therefore leaves the majority of query windows unmatched, resorting
to a single fallback value. SW-LTS assigns a non-zero kernel weight
to \emph{all} 19 entries for every query, with the nearest entries
contributing most. The selected bandwidth $\sigma^*=0.03$ (3\% of the
normalized range per dimension) enables genuine cross-rule interpolation:
a query window that differs from an observed pattern by one hedge
application still receives a substantial weight from that rule. The
result is MSE = 14,420, compared to 15,077 for HO-LTS---a 4.4\%
reduction.

\textbf{Matching behavior on larger datasets.}
For TAIEX-1999, Sunspot, and Temperature, coverage is 2.6--4.1\%, and
the exact-match mechanism already resolves most queries correctly.
In this regime, the kernel average over non-matching entries adds noise
rather than signal, and grid search selects $\sigma^*=0.01$---near
the exact-match limit. SW-LTS ($\lambda=3$) consequently matches
HO-LTS ($\lambda=3$) precisely on all four metrics for these datasets.
This graceful degradation is consistent with Proposition~1: SW-LTS
becomes LTS as $\sigma \to 0$.

\textbf{Comparison with traditional baselines.}
SW-LTS ($\lambda=3$) reduces MSE over Song \& Chissom (1993) by 98.2\%
on Alabama and 85.8\% on TAIEX-1999, and over Chen (1996) by 96.5\%
and 65.4\% respectively. These improvements reflect the cumulative
benefit of three design choices: the algebraically-grounded vocabulary
(vs.\ equal-width intervals), high-order rule extraction, and
kernel-weighted aggregation.

\textbf{LTS-PSO under default parameters.}
LTS-PSO is evaluated here with the same fixed $(\theta=0.57,
\alpha=0.49)$ as LTS, since PSO optimization would introduce
dataset-specific tuning. Its modified forecasting formula yields
MSE = 140,525 on Alabama, which is higher than standard LTS
(123,525); the formula change alone does not address the sparse-rule
problem and may introduce additional approximation error when rule
coverage is very low. On TAIEX-1999 (15,477) and Temperature (0.77),
LTS-PSO with optimized parameters would be expected to perform better.

\textbf{Scale invariance of $\sigma$.}
The optimal bandwidths across all datasets fall in $\sigma^* \in
[0.01, 0.03]$, a narrow range in the normalized $[0,1]^\lambda$ space,
despite the raw data spanning ranges from $[0, 296]$ (Sunspot) to
$[5{,}201, 9{,}039]$ (TAIEX). This confirms that normalizing the
kernel distance by the universe span is practically effective: a
single grid search strategy applies uniformly to all datasets.

\textbf{Computational cost.}
SW-LTS replaces an $O(1)$ hash lookup with an $O(M)$ kernel-weighted
sum, where $M = |\mathcal{R}| \leq T-\lambda$. Since $M$ is at most
a few hundred entries for the evaluated datasets, the added cost is
negligible relative to training time.

% ─────────────────────────────────────────────────────────────────────────────
\section{Conclusion}
\label{sec:conclusion}

We presented \textbf{SW-LTS}, a kernel-weighted extension of
Linguistic Time Series forecasting that addresses the sparse-rule
problem inherent in all existing LTS variants. By replacing binary
exact-match lookup with a Gaussian kernel aggregation over the full
LLRG in normalized HA semantic space, SW-LTS (i) exploits semantic
similarity across all rules, (ii) eliminates the sharp
matched/unmatched dichotomy, and (iii) introduces a scale-invariant
bandwidth $\sigma$ that admits consistent tuning across datasets.
We proved that SW-LTS strictly generalizes both LTS ($\sigma \to 0$)
and the HO-LTS global-mean fallback ($\sigma \to \infty$). Experiments
on four benchmarks confirm that the kernel mechanism provides
measurable accuracy gains when rule coverage is sparse (Alabama:
$-4.4\%$ MSE over HO-LTS) while converging to exact-match behavior
when coverage is adequate.

Future directions include: (1)~\textbf{adaptive bandwidth}---assigning
different $\sigma$ values to regions of the LLRG with heterogeneous
rule density; (2)~\textbf{multi-step prediction}---extending SW-LTS to
$h$-step-ahead forecasting via recursive or direct strategies;
(3)~\textbf{interval-valued output}---propagating kernel uncertainty
to produce forecast intervals $[\hat{y}^-_t, \hat{y}^+_t]$, exploiting
the natural interval semantics of HA; and (4)~\textbf{theoretical
analysis}---deriving approximation error bounds as a function of
$\sigma$, $M$, and the HA metric structure.

% ─────────────────────────────────────────────────────────────────────────────
\begin{thebibliography}{00}

\bibitem{song1993}
Q. Song and B. S. Chissom,
``Forecasting enrollments with fuzzy time series---Part~I,''
\textit{Fuzzy Sets and Systems}, vol.~54, no.~1, pp.~1--9, 1993.

\bibitem{song1994}
Q. Song and B. S. Chissom,
``Forecasting enrollments with fuzzy time series---Part~II,''
\textit{Fuzzy Sets and Systems}, vol.~62, no.~1, pp.~1--8, 1994.

\bibitem{chen1996}
S.-M. Chen,
``Forecasting enrollments based on fuzzy time series,''
\textit{Fuzzy Sets and Systems}, vol.~81, no.~3, pp.~311--319, 1996.

\bibitem{yu2005}
H.-K. Yu,
``Weighted fuzzy time series models for TAIEX forecasting,''
\textit{Physica A: Statistical Mechanics and its Applications},
vol.~349, no.~3--4, pp.~609--624, 2005.

\bibitem{catho2002ha}
N. C. Ho and W. Wechler,
``Hedge algebras: An algebraic approach to structure of sets of
linguistic truth values,''
\textit{Fuzzy Sets and Systems}, vol.~35, no.~3, pp.~281--293, 1990.

\bibitem{catho2002sqm}
N. C. Ho and N. V. Long,
``Fuzziness measure on complete hedge algebras and quantifying
semantics of terms in linear hedge algebras,''
\textit{Fuzzy Sets and Systems}, vol.~158, no.~4, pp.~452--471, 2007.

\bibitem{hieu2020lts}
N. D. Hieu, N. C. Ho, and V. N. Lan,
``Enrollment forecasting based on linguistic time series,''
\textit{Journal of Computer Science and Cybernetics}, vol.~36, no.~2,
pp.~119--137, 2020.

\bibitem{hieu2021holts}
N. D. Hieu,
``Scalable human knowledge about numeric time series variation and
its role in improving forecasting results,''
\textit{Journal of Computer Science and Cybernetics}, vol.~38, no.~2,
pp.~97--120, 2022.

\bibitem{hieu2022ltspso}
N. D. Hieu, M. V. Linh, and P. D. Phong,
``Linguistic time series forecasting based on particle swarm
optimization for hedge algebra parameter tuning,''
\textit{Journal of Computer Science and Cybernetics}, vol.~38, no.~4,
pp.~317--335, 2022.

\bibitem{hieu2023colts}
N. D. Hieu, M. V. Linh, and P. D. Phong,
``A co-optimization algorithm utilizing particle swarm optimization
for linguistic time series,''
\textit{Mathematics}, vol.~11, no.~7, p.~1597, 2023.
[Online]. Available: \url{https://doi.org/10.3390/math11071597}

\bibitem{kennedy1995pso}
J. Kennedy and R. Eberhart,
``Particle swarm optimization,''
in \textit{Proc. IEEE Int. Conf. Neural Networks (ICNN)},
Perth, Australia, Nov.--Dec. 1995, pp.~1942--1948.

\bibitem{zhan2024dfcnn}
T. Zhan, Y. He, Y. Deng, and Z. Li,
``Differential convolutional fuzzy time series forecasting,''
\textit{IEEE Transactions on Fuzzy Systems}, vol.~32, no.~2,
pp.~294--306, 2024.

\bibitem{wang2025bn}
B. Wang and X. Liu,
``Adaptive non-stationary fuzzy time series forecasting with
Bayesian networks,''
\textit{Sensors}, vol.~25, no.~5, p.~1598, 2025.
[Online]. Available: \url{https://doi.org/10.3390/s25051598}

\end{thebibliography}

\end{document}
